\documentclass[12pt,a4paper]{report}

% ---------- Encoding & Language ----------
\usepackage[T1]{fontenc}
\usepackage[utf8]{inputenc}
\usepackage[english]{babel}

% ---------- Page ----------
\usepackage[a4paper,margin=2.5cm]{geometry}
\usepackage{setspace}
\onehalfspacing
\usepackage{parskip}

% ---------- Useful packages ----------
\usepackage{graphicx}
\usepackage{float}
\usepackage{booktabs}
\usepackage{array}
\usepackage{longtable}
\usepackage{tabularx}
\usepackage{xcolor}
\usepackage{hyperref}
\usepackage{enumitem}
\usepackage{amsmath}
\usepackage{fancyhdr}
\usepackage{titlesec}
\usepackage{caption}
\usepackage{ifthen}
\usepackage{pdflscape}
\usepackage{pgfgantt}
\usepackage{tikz}

% ---------- Hyperref ----------
\hypersetup{
  colorlinks=true,
  linkcolor=blue!60!black,
  filecolor=magenta,
  urlcolor=cyan!60!black,
  pdftitle={PFE Report - Cloud VM Management Platform},
  pdfauthor={Your Name},
  pdfsubject={Final Year Project Report},
  pdfkeywords={PFE, OpenNebula, Next.js, NestJS, VM, Agile, UML}
}

% ---------- Header/Footer ----------
\pagestyle{fancy}
\fancyhf{}
\fancyhead[L]{PFE Report}
\fancyhead[R]{\leftmark}
\fancyfoot[C]{\thepage}

% ---------- Styling ----------
\titleformat{\chapter}[display]
  {\normalfont\bfseries\Huge}
  {\chaptertitlename\ \thechapter}{20pt}{\Huge}

% ---------- Custom commands ----------
\newcommand{\projectTitle}{Design and Implementation of a Multi-Service Cloud VM Management Platform}
\newcommand{\studentName}{Your Full Name}
\newcommand{\teammateName}{Teammate Full Name}
\newcommand{\supervisorName}{Supervisor Name}
\newcommand{\academicYear}{2025--2026}

% Safe figure include (compiles even if image not yet added)
\newcommand{\safeimage}[4][0.9\textwidth]{%
\begin{figure}[H]
\centering
\IfFileExists{#2}{
    \includegraphics[width=#1]{#2}
}{
    \fbox{\parbox{0.85\textwidth}{
    \centering
    \vspace{0.5em}
    Placeholder: add image file \texttt{#2}\\
    Description: #3
    \vspace{0.5em}
    }}
}
\caption{#3}
\label{#4}
\end{figure}
}

% Teammate-owned section placeholder
\newcommand{\serverowned}[1]{%
\begin{center}
\fcolorbox{red!60!black}{red!5}{
\parbox{0.92\textwidth}{
\textbf{[TO BE COMPLETED BY RSI/SERVER TEAMMATE]}\\
#1
}}
\end{center}
}

\begin{document}

% ---------- Cover ----------
\begin{titlepage}
    \centering
    \vspace*{2cm}
    {\Huge \textbf{\projectTitle} \par}
    \vspace{1.5cm}
    {\Large Final Year Project (PFE) Report\par}
    \vspace{2cm}
    \begin{tabular}{rl}
        \textbf{Student (Development):} & \studentName \\
        \textbf{Teammate (RSI/Server):} & \teammateName \\
        \textbf{Supervisor:} & \supervisorName \\
        \textbf{Academic Year:} & \academicYear \\
    \end{tabular}
    \vfill
    {\large Department of Computer Science\\
    University Name}
\end{titlepage}

% ---------- Front matter ----------
\pagenumbering{roman}
\chapter*{Acknowledgments}
Write a short acknowledgment (supervisor, institution, teammate, family).

\chapter*{Abstract}
This report presents the design and implementation of a cloud VM management platform based on a microservice architecture.
The solution includes a modern frontend, API gateway, authentication and user services, VM lifecycle service, and a worker layer.
The project follows Agile methodology with iterative sprint delivery and continuous integration practices.
Key outcomes include improved maintainability, scalable architecture, and a real-time terminal workflow for VM interaction.
Infrastructure and OpenNebula server operations are handled collaboratively with a dedicated RSI/server teammate.

\chapter*{Résumé (French - optional if required)}
(Optional) Add French summary if your institution requires it.

\tableofcontents
\listoffigures
\listoftables
\clearpage

\pagenumbering{arabic}

% ==========================
\chapter{General Introduction}
\section{Problem Statement}
Organizations, laboratories, and student innovation teams increasingly depend on virtualized infrastructure to deliver practical software projects under constrained time and hardware resources.
However, available VM administration solutions are frequently either too complex for daily academic use or too limited for modern development workflows.
In many cases, users must switch between multiple interfaces (authentication panel, VM dashboard, SSH client, and monitoring pages), which reduces productivity and increases operational errors.

The core problem addressed in this project is therefore the absence of a unified, developer-centered platform that combines VM lifecycle management, secure access, and interactive terminal capabilities in a single web experience.
The project also targets the pedagogical objective of building a production-inspired architecture using microservices, containerization, and iterative Agile delivery.

\section{Project Objectives}
\begin{itemize}[leftmargin=*]
    \item Build a full-stack platform for VM management and monitoring.
    \item Provide secure authentication and role-aware access control.
    \item Offer real-time terminal interaction for VM administration.
    \item Integrate with virtualization infrastructure (OpenNebula).
    \item Apply Agile methodology with sprint-based delivery.
\end{itemize}

From an engineering perspective, the project seeks to optimize three dimensions simultaneously: reliability, maintainability, and usability.
Reliability is pursued through service separation and explicit error handling; maintainability through modular service boundaries and typed contracts; and usability through a responsive dashboard and streamlined terminal experience.
The resulting platform is intended to remain extensible for future capabilities such as advanced observability, quota policies, and automated test pipelines.

\section{Report Organization}
This report is organized as follows.
Chapter 2 presents the context and requirement analysis.
Chapter 3 details the Agile project management framework, including backlog governance and sprint execution.
Chapter 4 introduces the analysis and design artifacts, especially UML views.
Chapter 5 explains implementation choices and technical realization.
Chapter 6 reserves infrastructure-specific work to the RSI/server teammate and formalizes integration boundaries.
Chapter 7 introduces the validation strategy while intentionally preserving pending test work.
Finally, Chapter 8 discusses results, limitations, and future extensions.

% ==========================
\chapter{Context and Requirement Analysis}
\section{Host Organization and Context}
The project was conducted in an academic final-year context where functional delivery quality and engineering rigor are both evaluated.
The team is composed of two contributors with complementary responsibilities: an application-development contributor (frontend/backend platform delivery) and an RSI/server contributor (infrastructure and OpenNebula operations).
This split enables clearer accountability but introduces coordination dependencies, particularly for deployment-ready validation scenarios.

Operationally, the project must satisfy both pedagogical constraints (documentation depth, UML completeness, reproducible development process) and practical constraints (limited timeline, infrastructure access windows, and progressive integration with server-side resources).

\section{Functional Requirements}
\begin{longtable}{p{0.2\textwidth} p{0.72\textwidth}}
\toprule
\textbf{ID} & \textbf{Requirement} \\
\midrule
FR1 & User registration and login (JWT-based authentication). \\
FR2 & User profile and SSH key management. \\
FR3 & VM provisioning, listing, and lifecycle actions (start/stop/restart/delete). \\
FR4 & Access real-time VM terminal from web interface. \\
FR5 & Role-aware dashboard and administration capabilities. \\
\bottomrule
\end{longtable}

\section{Non-Functional Requirements}
\begin{itemize}[leftmargin=*]
    \item Security: secure auth, controlled CORS, validated inputs.
    \item Performance: low-latency terminal streaming and responsive UI.
    \item Scalability: microservices and worker separation.
    \item Maintainability: modular codebase and clear service boundaries.
    \item Reliability: robust error handling and service health checks.
\end{itemize}

For terminal responsiveness, the qualitative target is an interaction feel close to local shell usage, with smooth command echo and continuous output streaming.
For maintainability, each service boundary should remain explicit and versionable, allowing incremental changes without full-system rewrites.
For operational robustness, build reproducibility through containerized execution is considered mandatory.

\section{Constraints and Assumptions}
\begin{itemize}[leftmargin=*]
    \item Infrastructure resources are managed by RSI/server side.
    \item Integration with OpenNebula depends on server-level configuration.
    \item Some advanced testing remains future work.
\end{itemize}

An important project assumption is that infrastructure APIs remain available and consistently configured across development and demonstration windows.
Any deviation in OpenNebula templates, network exposure, or access policy can directly affect functional demonstrations, which is why infrastructure-related sections are explicitly delegated and documented in a separate teammate handoff.

% ==========================
\chapter{Methodology and Project Management (Agile/Scrum)}
\section{Why Agile}
Agile was selected to support iterative value delivery, rapid feedback, and progressive risk reduction.

Given the progressive discovery nature of infrastructure integration, a sequential model (waterfall-style) would have increased late-stage risk.
Agile practices allowed the team to prioritize end-user value first (authentication, user operations, VM actions), then stabilize complex integration flows (real-time terminal, infrastructure dependencies).
Each sprint review was used as a decision point for backlog reprioritization.

\section{Roles and Responsibilities}
\begin{table}[H]
\centering
\begin{tabular}{p{0.25\textwidth} p{0.65\textwidth}}
\toprule
\textbf{Role} & \textbf{Responsibilities} \\
\midrule
Developer (You) & Frontend and backend service implementation, integration, API, UX, report writing. \\
RSI/Server Teammate & Server provisioning, OpenNebula setup, network/security hardening at infra level, monitoring setup. \\
Supervisor & Project guidance, milestone validation, quality oversight. \\
\bottomrule
\end{tabular}
\caption{Team roles and responsibilities}
\end{table}

\section{Product Backlog}
\begin{table}[H]
\centering
\small
\begin{tabular}{p{0.08\textwidth} p{0.45\textwidth} p{0.15\textwidth} p{0.12\textwidth} p{0.12\textwidth}}
\toprule
\textbf{ID} & \textbf{User Story} & \textbf{Priority} & \textbf{Estimation} & \textbf{Status} \\
\midrule
PB1 & As a user, I want to register/login securely. & High & 5 SP & Done \\
PB2 & As a user, I want to manage my profile and SSH keys. & High & 8 SP & Done \\
PB3 & As a user, I want to create and control VMs. & High & 13 SP & Done \\
PB4 & As a user, I want a real-time terminal for VM access. & High & 13 SP & In Progress \\
PB5 & As admin, I want system visibility and health checks. & Medium & 8 SP & In Progress \\
PB6 & As team, we want production-ready server integration. & High & 13 SP & RSI Side \\
\bottomrule
\end{tabular}
\caption{Initial product backlog}
\end{table}

Backlog refinement was performed at least once per sprint.
Stories were estimated using story points and reassessed whenever technical uncertainty changed.
High-risk stories (terminal synchronization and infrastructure integration) were decomposed into smaller implementation tasks to reduce delivery uncertainty and improve review clarity.

\section{Sprint Planning and Execution}
\subsection{Sprint Breakdown}
\begin{itemize}
    \item Sprint 1: Core architecture and authentication.
    \item Sprint 2: User/SSH key services and dashboard basics.
    \item Sprint 3: VM service integration and lifecycle actions.
    \item Sprint 4: Real-time terminal stabilization and hardening.
    \item Sprint 5: Documentation, validation, and finalization.
\end{itemize}

Each sprint followed the same control cycle: planning (scope commitment), daily synchronization (impediment detection), review (demonstrable increments), and retrospective (process improvement actions).
This closed feedback loop made it possible to continuously improve both technical quality and team collaboration practices.

\subsection{Sprint Gantt (Template)}
\begin{figure}[H]
\centering
\begin{ganttchart}[
    vgrid, hgrid,
    x unit=0.7cm,
    y unit title=0.5cm,
    y unit chart=0.6cm,
    title label font=\bfseries\small,
    bar label font=\small
]{1}{15}
\gantttitle{Project Timeline (Weeks)}{15} \\
\gantttitlelist{1,...,15}{1} \\
\ganttbar{Sprint 1}{1}{3} \\
\ganttbar{Sprint 2}{4}{6} \\
\ganttbar{Sprint 3}{7}{9} \\
\ganttbar{Sprint 4}{10}{12} \\
\ganttbar{Sprint 5}{13}{15}
\end{ganttchart}
\caption{Sprint timeline (update with your actual dates)}
\end{figure}

\section{Sprint Ceremonies}
\textbf{Sprint Planning:} prioritized backlog items were transformed into executable tasks with explicit acceptance criteria.

\textbf{Daily Synchronization:} blockers were captured early, especially those linked to service integration and environment availability.

\textbf{Sprint Review:} each increment was demonstrated through working software artifacts (UI behavior, API actions, VM operations).

\textbf{Retrospective:} process improvements were recorded and applied in the subsequent sprint (e.g., earlier interface contract checks, better risk visibility, and tighter integration checklists).

\section{Sprint and Backlog Metrics}
\begin{table}[H]
\centering
\begin{tabular}{p{0.25\textwidth} p{0.65\textwidth}}
\toprule
\textbf{Metric} & \textbf{Interpretation} \\
\midrule
Planned vs Delivered Story Points & Measures sprint predictability and commitment quality. \\
Carry-over Stories & Indicates estimation accuracy and dependency pressure. \\
Blocked Tasks per Sprint & Highlights external or infrastructure bottlenecks. \\
Defect Notes per Increment & Tracks quality trends and technical debt signals. \\
\bottomrule
\end{tabular}
\caption{Agile monitoring indicators used in this project}
\end{table}

% ==========================
\chapter{System Analysis and Design}
\section{Global Vision}
\safeimage{figures/architecture-overview.png}{Global architecture overview}{fig:global-arch}

The platform is organized around a gateway-centric interaction model.
Frontend clients communicate with the gateway, which routes requests to specialized backend services.
This design improves separation of concerns while preserving a unified entry point for authentication, authorization, and observability.

\section{UML Use Case Diagram}
\safeimage{figures/uml-usecase.png}{Use case diagram for platform actors and interactions}{fig:usecase}

\section{Domain/Class Diagram}
\safeimage{figures/uml-class.png}{Class/domain model diagram}{fig:class}

\section{Component Diagram}
\safeimage{figures/uml-component.png}{Component diagram (frontend, gateway, services, worker, DB)}{fig:component}

\section{Sequence Diagrams}
\safeimage{figures/uml-seq-auth.png}{Authentication flow sequence diagram}{fig:seq-auth}
\safeimage{figures/uml-seq-vm-lifecycle.png}{VM lifecycle sequence diagram}{fig:seq-vm}
\safeimage{figures/uml-seq-terminal.png}{Real-time terminal interaction sequence diagram}{fig:seq-terminal}

\section{Deployment Diagram}
\safeimage{figures/uml-deployment.png}{Deployment diagram with infrastructure nodes}{fig:deployment}
\serverowned{Add precise server topology: OpenNebula hosts, hypervisors, network segmentation, firewall placement, reverse proxy, TLS termination, and monitoring agents.}

The deployment view distinguishes application-plane components (frontend and microservices) from infrastructure-plane components (OpenNebula and host-level controls).
This distinction is fundamental to clarify ownership boundaries between development and RSI/server activities.

% ==========================
\chapter{Implementation}
\section{Technology Stack}
\begin{table}[H]
\centering
\begin{tabular}{p{0.3\textwidth} p{0.6\textwidth}}
\toprule
\textbf{Layer} & \textbf{Technology} \\
\midrule
Frontend & Next.js + TypeScript + Tailwind CSS \\
Backend API & NestJS microservices (Gateway, Auth, User, VM) \\
Data & Prisma + relational database \\
Async/Workers & Python worker modules \\
Deployment & Docker + Docker Compose \\
Virtualization Integration & OpenNebula (RSI/server-managed) \\
\bottomrule
\end{tabular}
\caption{Technology stack}
\end{table}

The stack was selected to maximize developer productivity while preserving production-like patterns.
Next.js and TypeScript support robust frontend iteration, while NestJS enables modular backend services with clear architectural conventions.
Prisma contributes schema-driven consistency between domain model and persistence.

\section{Frontend Highlights}
The frontend follows a route-structured design with dedicated areas for authentication, dashboard operations, VM management, profile settings, and SSH key administration.
Particular attention was given to terminal interaction ergonomics (clear output rendering, command flow continuity, and user feedback on connection state).
The UI was implemented to keep operational actions discoverable and reduce context switching.

\section{Backend Highlights}
\begin{itemize}
    \item Gateway service for routing/proxy.
    \item Auth service with JWT and guards.
    \item User service for profile and SSH keys.
    \item VM service for lifecycle management and terminal channeling.
    \item Worker service for automation/background updates.
\end{itemize}

At implementation level, service interfaces were intentionally kept explicit to reduce coupling.
Authentication and authorization concerns were isolated from VM orchestration logic.
This boundary discipline simplifies future evolution, including policy enrichment and audit extension.

\section{Data Model Alignment}
The domain model is centered around the \texttt{User} entity and extends to \texttt{SshKey}, \texttt{VirtualMachine}, \texttt{Plan}, \texttt{UserQuota}, \texttt{RefreshToken}, \texttt{Payment}, and \texttt{Notification}.
Entity relations were designed to support both operational workflows (provisioning, terminal access) and account-level governance (quota and billing evolution).

\section{Security and Hardening}
\begin{itemize}
    \item Input validation and DTO constraints.
    \item Auth guards and token checks.
    \item Error handling strategy and secure defaults.
    \item CORS and environment variable governance.
\end{itemize}

Hardening activities focused on practical risk reduction rather than theoretical checklists.
The selected controls prioritize secure defaults, constrained request surfaces, and failure behavior clarity.
This approach keeps the platform evolvable while improving baseline resilience.

% ==========================
\chapter{Infrastructure and OpenNebula Integration}
\section{Integration Objectives}
Describe what the platform expects from infrastructure APIs and services.

\section{Server Architecture}
\serverowned{Provide full server architecture details:
\begin{itemize}
    \item OpenNebula cluster topology.
    \item VM template lifecycle and image pipeline.
    \item Network/VLAN mapping and access model.
    \item Storage configuration and backup policy.
\end{itemize}
}

\section{Provisioning and Operations}
\serverowned{Explain deployment scripts, systemd/services, reverse proxy configs, SSL certificates, secrets management, and recovery procedures.}

\section{Monitoring and Incident Handling}
\serverowned{Document logs, metrics, alert rules, incident response process, and observed issues/resolutions.}

% ==========================
\chapter{Validation and Testing}
\textbf{Important note:} This chapter is intentionally incomplete and will be finalized later.

\section{Current Validation Status}
At this stage, validation mainly covers implemented build and lint workflows for relevant project modules, plus functional manual checks for key user journeys.
Comprehensive automated and non-functional campaigns are intentionally postponed and will be completed in a subsequent iteration.

\section{Planned Test Strategy (TO DO)}
\begin{itemize}
    \item Unit tests by service.
    \item API integration tests.
    \item Terminal real-time behavior tests.
    \item Security and penetration checks.
    \item Performance benchmarks.
\end{itemize}

The final strategy will include traceability from requirement IDs to test cases, with explicit pass/fail evidence and defect references.
For terminal validation, special attention will be given to stream continuity, reconnection handling, and command latency perception.

\section{Test Cases Table (Template)}
\begin{longtable}{p{0.12\textwidth} p{0.45\textwidth} p{0.15\textwidth} p{0.18\textwidth}}
\toprule
\textbf{ID} & \textbf{Scenario} & \textbf{Expected} & \textbf{Result} \\
\midrule
T1 & Login with valid credentials & JWT issued & TO DO \\
T2 & Add SSH key & Key stored safely & TO DO \\
T3 & Create VM & VM provisioned & TO DO \\
T4 & Open terminal session & Interactive shell works & TO DO \\
\bottomrule
\end{longtable}

% ==========================
\chapter{Discussion and Future Work}
\section{Achievements}
The project successfully delivered a coherent multi-service platform with an integrated user workflow for authentication, account management, VM lifecycle actions, and terminal interaction.
Architecturally, the separation into dedicated services improved clarity of responsibilities and supports long-term extension.

\section{Limitations}
Current limitations include incomplete automated testing coverage and dependence on infrastructure readiness for full end-to-end production-like validation.
Additionally, advanced observability and policy automation remain partially defined and require dedicated future iterations.

\section{Future Improvements}
\begin{itemize}
    \item Full automated test pipeline.
    \item Enhanced observability dashboard.
    \item Multi-tenant policy and quota management.
    \item AI-assisted incident diagnostics.
\end{itemize}

% ==========================
\chapter{General Conclusion}
This project delivered a modular and scalable VM management platform aligned with Agile practices.
The implemented architecture separates concerns across services and supports progressive hardening.
Further work focuses on complete test coverage and deeper production infrastructure optimization.

Beyond feature delivery, the project demonstrates a structured engineering process where architecture, implementation, and project management artifacts remain consistent.
This coherence is central to PFE evaluation because it links technical outputs to methodical execution and measurable progress.

% ---------- References ----------
\begin{thebibliography}{99}
\bibitem{nextjs} Next.js Documentation. \url{https://nextjs.org/docs}
\bibitem{nestjs} NestJS Documentation. \url{https://docs.nestjs.com}
\bibitem{prisma} Prisma Documentation. \url{https://www.prisma.io/docs}
\bibitem{docker} Docker Documentation. \url{https://docs.docker.com}
\bibitem{opennebula} OpenNebula Documentation. \url{https://docs.opennebula.io}
\bibitem{scrumguide} The Scrum Guide. \url{https://scrumguides.org}
\end{thebibliography}

\appendix
\chapter{Appendix A: Sprint Retrospective Templates}
\section*{What went well}
\section*{What did not go well}
\section*{Action items}

\chapter{Appendix B: API Endpoints (Template)}
List your most important endpoints by service.

\end{document}
